% Cover illustration: observations, a learned multimodal density, and samples.
\documentclass[border=2mm]{standalone}
\usepackage{amsmath,amssymb}
% Chinese labels in figures; fonts are shared with the book.
\usepackage[UTF8,fontset=none]{ctex}
% CJK fonts for the book, the figures and the slides.
% Noto Sans CJK SC (SIL Open Font License) is the same typeface as Source Han Sans SC.
% macOS: brew install --cask font-noto-sans-cjk-sc
% Debian/Ubuntu: apt install fonts-noto-cjk
\setCJKmainfont[BoldFont={Noto Sans CJK SC Bold}]{Noto Sans CJK SC}
\setCJKsansfont[BoldFont={Noto Sans CJK SC Bold}]{Noto Sans CJK SC}
\setCJKmonofont{Noto Sans CJK SC}


% Shared visual language for the figures of the book.
% Every figure in this directory loads this file so the chapter reads as one system.
% NOTE: no blank lines are allowed inside \tikzset{...} (a blank line emits \par
% and breaks key scanning).
\NeedsTeXFormat{LaTeX2e}
\ProvidesPackage{figstyle}[2026/10/21 Figure style for Free Energy to LLMs]

\RequirePackage{tikz}
\usetikzlibrary{arrows.meta,positioning,calc,fit,shapes.geometric,shapes.misc,%
  backgrounds,matrix,decorations.pathreplacing,decorations.markings,patterns.meta}

% ---------------------------------------------------------------- palette
% Same values as book/preamble/style.tex, prefixed with F.
\definecolor{FNavy}{HTML}{14324A}
\definecolor{FTeal}{HTML}{0E7490}
\definecolor{FSky}{HTML}{E7F5F8}
\definecolor{FOrange}{HTML}{D97706}
\definecolor{FAmber}{HTML}{FFF3D6}
\definecolor{FInk}{HTML}{1F2937}
\definecolor{FSlate}{HTML}{64748B}
\definecolor{FMist}{HTML}{F5F7FA}
\definecolor{FRose}{HTML}{9F1239}
\definecolor{FRoseMist}{HTML}{FFF1F2}

% ---------------------------------------------------------------- semantics
% Teal   = probability / model / data object
% Orange = training objective, cost, or the quantity moved through the system
% Navy   = neutral structure
% Slate  = secondary label or metadata
% Rose   = failure mode, hazard, masked-out region
\tikzset{
  % Lay figures out in millimetres.  A figure designed 168mm wide is included at
  % \textwidth, so its fonts appear at exactly the size written here.
  every picture/.style={x=1mm,y=1mm,font=\small,line cap=round,line join=round},
  % ---- block styles
  nbase/.style={draw,rounded corners=1.2pt,align=center,inner sep=2.4pt,%
    minimum height=7mm,text=FInk},
  nteal/.style={nbase,draw=FTeal,fill=FSky},
  norange/.style={nbase,draw=FOrange,fill=FAmber},
  nnavy/.style={nbase,draw=FNavy,fill=white},
  nsoft/.style={nbase,draw=FSlate!55,fill=FMist},
  nrose/.style={nbase,draw=FRose,fill=FRoseMist},
  ghost/.style={draw=none,fill=none,align=center,inner sep=1pt},
  plain/.style={draw=none,fill=none,align=center,inner sep=0pt},
  % ---- arrows
  ar/.style={-{Latex[length=1.7mm,width=1.35mm]},draw=FSlate,line width=0.75pt},
  ard/.style={ar,dashed},
  arteal/.style={ar,draw=FTeal},
  arorange/.style={ar,draw=FOrange},
  arrose/.style={ar,draw=FRose},
  thickar/.style={ar,line width=1.2pt},
  % ---- text roles
  ftitle/.style={font=\large\bfseries,text=FNavy,align=center},
  ptitle/.style={font=\small\bfseries,text=FNavy,align=center},
  ptitleo/.style={font=\small\bfseries,text=FOrange,align=center},
  lbl/.style={font=\scriptsize,text=FSlate,align=center,inner sep=1pt},
  lblt/.style={lbl,text=FTeal},
  lblo/.style={lbl,text=FOrange},
  lbln/.style={lbl,text=FNavy},
  lblr/.style={lbl,text=FRose},
  sub/.style={font=\scriptsize,text=FSlate!90,align=center,inner sep=1pt},
  monob/.style={font=\ttfamily\scriptsize,text=FInk,inner sep=1pt},
  mathb/.style={font=\small,text=FInk,inner sep=1pt},
  eqn/.style={font=\scriptsize,inner sep=1.5pt,text=FInk},
  % ---- decorations
  brace/.style={decorate,decoration={brace,amplitude=2.4pt},draw=FSlate,%
    line width=0.6pt},
  bracet/.style={brace,draw=FTeal},
  braceo/.style={brace,draw=FOrange},
  axis/.style={draw=FSlate!55,line width=0.6pt},
  guide/.style={draw=FSlate!45,line width=0.5pt,dashed},
  halolabel/.style={font=\scriptsize,text=FInk,fill=white,fill opacity=0.85,%
    text opacity=1,inner sep=1.4pt,align=center},
}

% Figure title above the drawing, anchored at the figure's top-left.
\newcommand{\figtitle}[2][0]{%
  \node[ftitle,anchor=north west,outer sep=0pt] at (#1,0) {#2};}

% Rounded panel behind a subgroup of nodes.
\newcommand{\figpanel}[3][]{%
  \node[draw=FSlate!35,fill=#3,rounded corners=2pt,inner sep=2.6mm,#1] {#2};}

\endinput

\begin{document}
\begin{tikzpicture}[x=1mm,y=1mm]
  % Three quiet fields make the visual read as one transformation.
  \fill[FSky!75] (28,37) circle (21mm);
  \fill[FMist] (80,37) circle (22mm);
  \fill[FAmber!62] (132,37) circle (21mm);
  \draw[arteal,line width=1.1pt] (51,37)--(57,37);
  \draw[arorange,line width=1.1pt] (103,37)--(109,37);

  % Two clusters of observed data.
  \foreach \x/\y/\r in {
    16/44/1.0,19/48/1.25,22/43/1.1,24/49/1.45,27/45/1.1,
    30/50/1.0,31/42/1.2,35/46/1.0,37/50/1.3,
    18/28/1.05,21/23/1.3,24/29/1.0,27/25/1.2,30/31/1.0,
    32/25/1.45,35/29/1.1,39/24/1.0} {
    \fill[FTeal!90] (\x,\y) circle (\r mm);
  }
  \fill[FTeal!40] (13,39) circle (.7mm);
  \fill[FTeal!40] (40,39) circle (.7mm);

  % The model represents density around both modes rather than one mean.
  \foreach \a/\b/\opacity in {13/8/35,10/6/50,7/4/75} {
    \draw[FTeal!\opacity,line width=.75pt,rotate around={-22:(72,44)}]
      (72,44) ellipse (\a mm and \b mm);
  }
  \foreach \a/\b/\opacity in {13/8/35,10/6/50,7/4/75} {
    \draw[FOrange!\opacity,line width=.75pt,rotate around={23:(89,29)}]
      (89,29) ellipse (\a mm and \b mm);
  }
  \fill[FTeal] (72,44) circle (1.25mm);
  \fill[FOrange] (89,29) circle (1.25mm);

  % New samples follow the same two-mode geometry with different locations.
  \foreach \x/\y/\r in {
    120/45/1.1,123/49/1.0,126/43/1.25,128/48/1.0,131/45/1.45,
    135/51/1.05,138/44/1.2,142/47/1.0,
    119/27/1.0,122/23/1.3,125/30/1.0,128/25/1.15,132/28/1.0,
    135/22/1.4,138/29/1.0,142/25/1.1} {
    \fill[FOrange!95] (\x,\y) circle (\r mm);
  }
  \fill[FOrange!45] (116,39) circle (.7mm);
  \fill[FOrange!45] (145,36) circle (.7mm);

  \node[font=\sffamily\small,text=FNavy] at (28,7) {概率分布};
  \node[font=\sffamily\small,text=FNavy] at (80,7) {统计学习};
  \node[font=\sffamily\small,text=FNavy] at (132,7) {采样与推断};
  \node[font=\scriptsize,text=FSlate] at (28,0) {what to represent};
  \node[font=\scriptsize,text=FSlate] at (80,0) {how to fit};
  \node[font=\scriptsize,text=FSlate] at (132,0) {how to generate};
\end{tikzpicture}
\end{document}
